%% kapa: exam=1 round=30 year=2019 period=1 form=A subject=economics subject_ko=경제학원론 seq=25 number=65
%% source: Q-Net 감정평가사 기출문제 (cst00309) artlSeq=5205205 file=extracted/1008/5205205/2019년 제30회 감정평가사 1차시험 문제지/2019년 제30회 감정평가사 1차시험 문제지/2019년도 제30회 감정평가사 1차 1교시 A형.hwp
%% conversion: kordoc 4.24.51 -> tools/gen1.py (auto); review=proofread-sonnet; flags=none
\begin{question}{65}
개방경제 甲국의 국민소득 결정모형이 다음과 같다. 특정 정부지출 수준에서 경제가 균형을 이루고 있으며 정부도 균형예산을 달성하고 있을 때, 균형에서 민간저축은? (단, ${ Y }$는 국민소득, ${ C }$는 소비, ${ I }$는 투자, ${ G }$는 정부지출, ${ T }$는 조세, ${ X }$는 수출, ${ M }$은 수입이다.)
\begin{framedtext}
○ ${ Y=C+I+G+(X-M) }$ \par
○ ${ C=150+0.5(Y-T) }$ \par
○ ${ I=200 }$ \par
○ ${ T=0.2Y }$ \par
○ ${ X=100 }$ \par
○ ${ M=50 }$
\end{framedtext}
\begin{choices}
  \item 150
  \item 200
  \item 225
  \item 250
  \item 450
\end{choices}
\end{question}
\answer{4}
